\documentclass[10pt]{article}
\usepackage[T1]{fontenc}
\usepackage[utf8]{inputenc}
\usepackage[french]{babel}
\usepackage{amsmath,amssymb,mathrsfs,mathpazo}
\usepackage{geometry}
\geometry{a4paper,margin=2cm}
\usepackage{graphicx,array,tabularx,tabularray,colortbl,multirow,multicol,booktabs,enumitem,pifont,ulem,fancybox,diagbox,cancel}
\usepackage{pstricks,pstricks-add,pst-plot,pst-tree,pst-node,pst-text,pst-func,pst-math,pst-bspline,pst-3dplot,pst-eucl,pst-solides3d}
\let\pscalculate\undefined
\usepackage{pst-bezier}
\usepackage{tikz,pgfplots,tkz-tab}
\usetikzlibrary{arrows,arrows.meta,patterns,calc,positioning,trees,decorations.markings,shapes.geometric,shapes.misc,shapes.arrows,intersections,decorations.pathmorphing,decorations.pathreplacing,backgrounds,shadows,babel}
\pgfplotsset{compat=1.18}
\usepgfplotslibrary{fillbetween}
\usepackage[np]{numprint}
\usepackage{listings,makecell,pas-tableur,twoopt}
\let\barre\undefined
\usepackage[pstricks]{tablvar}
\definecolor{darkgreen}{RGB}{0,100,0}
\providecommand{\red}{\color{red}}
\providecommand{\blue}{\color{blue}}
\providecommand{\green}{\color{green}}
\providecommand{\black}{\color{black}}
\providecommand{\white}{\color{white}}
\providecommand{\gray}{\color{gray}}
\providecommand{\cyan}{\color{cyan}}
\providecommand{\magenta}{\color{magenta}}
\providecommand{\yellow}{\color{yellow}}
\providecommand{\violet}{\color{violet}}
\providecommand{\orange}{\color{orange}}
\providecommand{\euro}{\texteuro}
\providecommand{\EUR}{\texteuro}
\providecommand{\ieme}{\textsuperscript{e}}
\providecommand{\degres}{\ensuremath{^\circ}}
\providecommand{\degree}{\ensuremath{^\circ}}
\providecommand{\vv}[1]{\overrightarrow{#1}}
\providecommand{\Coord}[1]{#1}
\providecommand{\decosix}{}
\providecommand{\starredbullet}{\ensuremath{\bullet}}
\pagestyle{empty}
\setlength{\parindent}{0pt}
\usepackage{adjustbox}\DeclareUnicodeCharacter{2236}{:}\DeclareUnicodeCharacter{2190}{\ensuremath{\leftarrow}}
\providecommand{\R}{\mathbb{R}}
\providecommand{\N}{\mathbb{N}}
\providecommand{\Z}{\mathbb{Z}}
\providecommand{\Q}{\mathbb{Q}}
\providecommand{\D}{\mathbb{D}}
\providecommand{\C}{\mathbb{C}}
\providecommand{\e}{{\mathrm{e}}}
\providecommand{\vect}[1]{\overrightarrow{#1}}
\providecommand{\vectt}[1]{\overrightarrow{\mathrm{#1}}}
\providecommand{\vvt}[1]{\overrightarrow{\mathrm{#1}}}
\providecommand{\np}[2][]{{#2~#1}}
\providecommand{\vv}[1]{{\overrightarrow{#1}}}
\renewcommand{\d}{{\mathrm{d}}}
\providecommand{\degres}{{^\circ}}
\providecommand{\degre}{{^\circ}}
\providecommand{\blue}{\color{blue}}
\providecommand{\red}{\color{red}}
\providecommand{\enskip}{\quad}
\providecommand{\llbracket}{\lbrack\!\lbrack}
\providecommand{\rrbracket}{\rbrack\!\rbrack}
\providecommand{\euro}{{\text{€}}}
\providecommand{\no}{{\text{n°}}}
\providecommand{\interval}[2]{{[#1;#2]}}
\providecommand{\intervFO}[2]{{[#1;#2[}}
\providecommand{\dotfill}{\ldots}

\providecommand{\ds}{\displaystyle}
\providecommand{\textding}[1]{#1}
\providecommand{\Oij}{\(\mathrm{O};\vec{\imath},\vec{\jmath}\)}
\providecommand{\Oijk}{\(\mathrm{O};\vec{\imath},\vec{\jmath},\vec{k}\)}
\providecommand{\Ouv}{\(\mathrm{O};\vec{u},\vec{v}\)}
\providecommand{\intervFF}[2]{\ensuremath{[#1;#2]}}
\renewcommand{\euro}{\ifmmode\text{\texteuro}\else\texteuro\fi}
\DeclareUnicodeCharacter{00B0}{\ifmmode^{\circ}\else\textdegree\fi}
\DeclareUnicodeCharacter{20AC}{\ifmmode\text{\texteuro}\else\texteuro\fi}
% Énoncés : Éducation nationale. Transcription et corrigés : APMEP, Denis Vergès et contributeurs. Découpage et classement : CoopMaths, Sébastien Lozano, Éric Elter et contributeurs. Adaptation de la présentation : Maths Explorer.
% CC BY-SA : https://coopmaths.fr/www/about/
% https://coopmaths.fr/alea/static/eam/2026/tex/eam_specifique_2026_06_polynesie_automatismes_cor.tex
\begin{document}
\noindent\textit{APMEP / CoopMaths. CC BY-SA. Présentation adaptée par Maths Explorer.}\par\medskip
\subsubsection*{Question 1}

$\dfrac{1}{3} \times \dfrac{6}{5}-\dfrac{1}{5} =  \dfrac{1 \times 2 \times 3}{3 \times 5} - \dfrac{1}{5} = \dfrac{2}{5} - \dfrac{1}{5} = \dfrac{1}{5}$

Bonne réponse : \textbf{C.~~} $\dfrac{1}{5}$


\subsubsection*{Question 2}

$\begin{aligned}
	3 x+2= -5 x+4
	&\iff 3x + 5x + 2 = 4 \\
	&\iff  8x + 2 =  4  \\
	&\iff  8x  =  2 \\
	&\iff x = \dfrac{2}{8} = \dfrac{1}{4}
\end{aligned}$

Bonne réponse : \textbf{C.~~} $x= \dfrac{1}{4}$

\subsubsection*{Question 3}

$(2x - 3)(2x + 3) = 2x\times 2x + 2x\times 3 + (-3)\times 2x + (-3)\times 3 = 4x^2 - 9$

Bonne réponse : \textbf{D.~~}$4x^2 - 9$

\subsubsection*{Question 4}

$5~\% = \dfrac{5}{100} = \dfrac{20}{400}$. L'effectif total des animaux est 400.

Bonne réponse : \textbf{B.~~} 400

\subsubsection*{Question 5}
$\np[km]{30} = \np[h]{0.5} \times \np[km/h]{60}$ et $\np[h]{0.5} = \np[min]{30}$.

Bonne réponse : \textbf{A.~~} 30

\subsubsection*{Question 6}
Augmenter de $5~\%$ par an c'est multiplier par \np{1.05} chaque année. Au bout de trois ans cela revient à multiplier par $\np{1.05}\times\np{1.05}\times\np{1.05} = \np{1.05}^3$.

Bonne réponse : \textbf{B.~~} $1,05^{3}$


\subsubsection*{Question 7}
Par lecture graphique, pour $x = 0$, $y = -1$. Seules les réponses A et B peuvent convenir.

Par lecture graphique, pour $x = 4$, $y = 0$. Par calcul : $4\times 4 - 1 = 15$  et $\dfrac{1}{4}\times 4 - 1 = 0$.

Bonne réponse : \textbf{B.~~} $y=\dfrac{1}{4} x - 1$


\subsubsection*{Question 8}
Bonne réponse : \textbf{A.~~}$f$ est croissante sur $[-2 ; 0]$

\end{document}
